\thispagestyle{empty}
\jaliningfigures   % Arabic numerals in lining figures (the folia reset it)
\begin{center}
\vspace*{26mm}
{\fontsize{26}{30}\selectfont\addfontfeatures{LetterSpace=10}SIDEREVS NVNCIVS}\\[4mm]
{\fontsize{20}{24}\selectfont 星\,界\,の\,使\,者}\\[10mm]
{\large Galilaeus Galilaeus · ガリレオ・ガリレイ}\\[2mm]
{\small Venetiis, MDCX · ヴェネツィア、1610年}\\[20mm]
{\color{edgray}\small ラテン語・日本語対訳}\\[6mm]
{\small metachow\ 編訳}\\[1mm]
{\color{edgray}\footnotesize Claude（Anthropic）協力：翻訳・校閲・組版}
\end{center}

\vfill
\noindent\begin{minipage}{\textwidth}
\small\color{edgray}\setlength{\parindent}{1em}\linespread{1.25}\selectfont\sloppy
\textbf{凡例}\quad
ラテン語本文は A. Favaro 編『ガリレオ著作集』国家版第3巻（1892）を底本とし、1610年ヴェネツィア初版と1ページずつ照合した。初版はスミソニアン図書館所蔵本（Internet Archive によるデジタル化）による。綴りは現代化した。æ、œ は ae、oe と書き分け、長い s（ſ）は s とし、u と v は音価によって区別し、i はそのままとした（\textit{Iuppiter, Ianuarii}）。略語はすべて展開し、綴りは古典的な形に統一した（\textit{quattuor, numquam, intellego, exspecto}）。普通名詞は小文字とし、天体名（\textit{Sol, Luna, Tellus, Iuppiter}）は大文字を残した。

本書の1ページは初版の1ページに対応する。柱は初版の柱をそのまま写し、外側の角括弧内に初版の丁付を示した（[7r] は第7丁の表）。ラテン語は初版と同じ箇所でページを改め、ページをまたいで分かれた単語はその語の始まるページに置いた。日本語は語順がラテン語と異なるため文単位で配置し、ラテン語の文が始まるページに置いた。

書体は初版の使い分けに従った。初版のローマン体は正体、日本語は明朝体とし、イタリック体はイタリック、日本語は教科書体とした。初版が大文字で強調した語は字間を空けた大文字とし、日本語では傍点を付した。常用漢字外の字と特別な読みの語には、初出に限りルビを付した。

日本語訳はラテン語から直接訳した新訳である。挿図のうち、望遠鏡の光路図、幾何学的な図、星図、木星の衛星の図は初版を実測してベクター図として描き直した。月面の銅版画、木版の装飾頭文字、飾り模様と印刷者標章は初版の画像から取り、黄ばみを除き、折り目を修復した。

翻訳は AI の協力によるもので、ラテン語の専門家による校閲はまだ経ていない。誤りのご指摘を歓迎する。本書の訳文・注・組版は CC BY 4.0 で公開する。ラテン語原文と初版の画像はパブリックドメインである。ソース：github.com/metachow/sidereus-nuncius
\end{minipage}
\cleardoublepage
